\chapter{Conclusion and Future Perspectives}
\label{ch:conclusion}

\section{Conceptual and Methodological Synthesis}
The development of the NECTA (Neural Connectivity Time-resolved Analysis) framework represents a critical paradigm shift from traditional static graph representations to the continuous exploration of the neural "dynome." Historically, functional connectivity analyses in electrophysiology have collapsed the temporal dimension, effectively masking the highly volatile, millisecond-by-millisecond network state transitions that underlie sensory integration. By establishing a high-performance computational pipeline, this work successfully bridges the gap between raw high-density multi-electrode arrays and rigorous topological inference. 

Methodologically, NECTA directly addresses two core bottlenecks in modern neuroinformatics: the Big Data processing limitation and topological false positives. Through zero-phase filtering and the strategic application of robust phase-coupling metrics (such as iCoh and wPLI), the pipeline systematically mitigates the confounding effects of localized spatial crosstalk and volume conduction physical artifacts. Combined with memory-efficient array structures, the framework demonstrates that tracking time-resolved directional workflows is computationally viable without inducing catastrophic Out-Of-Memory (OOM) failures, thereby establishing a reproducible benchmark for large-scale cortical circuit mapping.

\section{Neurobiological Implications of the Dynome}
The empirical application of NECTA to high-density feline electrophysiological recordings yielded substantial insights into the functional architecture of the primary visual cortex. Central to these findings is the unmasking of the localized neurobiological role of the Transition Zone (TZ). Our statistical validation in Chapter 4 demonstrated that under a 15\% proportional density threshold in the Low Gamma band (30--59 Hz), the TZ operates with a remarkably high processing load, exhibiting a significantly elevated Node Strength ($M = 9.60$, $SD = 3.50$) compared to the major extrastriate Area 18 ($M = 7.25$, $SD = 3.11$; $p = 0.030$, effect size $r = 0.453$). 

Biologically, this strong effect size demonstrates that despite its restricted anatomical boundary compared to the massive retinotopic expanses of striate and extrastriate cortices, the TZ acts as an ultra-dense inter-hemispheric convergence hub. The high topological variance and causal bimodality captured by NECTA suggest that the TZ serves as a flexible, self-tuning modulator, dynamically routing Source/Sink communication channels depending on the immediate state of visual processing. Crucially, because these non-stationary coupling patterns were recorded under a controlled anesthetized animal model, they reflect fundamental processes of primary intrinsic cortical plasticity, operating entirely independent of top-down conscious attention or cognitive modulations.

\section{Technical Limitations}
An honest assessment of the NECTA framework highlights inherent operational boundaries that shape the interpretation of the dynome. First, a fundamental trade-off between spatial and temporal configurations persists; while multi-electrode electrophysiology yields microsecond-level temporal precision ideal for fast synchronization metrics, its spatial coverage remains anchored to specific, surgically restricted Regions of Interest (ROIs), lacking the holistic whole-brain macro-topology provided by lower-resolution metabolic methods such as fMRI. 

Second, the structural topology extracted by the pipeline remains sensitive to operational threshold choices. The enforcement of a specific proportional density limiar (e.g., 15\%) and the discretization of continuous time series into brief, arbitrary sliding windows operate as mathematical filters. While these constraints are mathematically required to isolate non-stationary patterns from stochastic noise, they inevitably establish a boundary condition that bounds the absolute visibility of the wider connectivity continuum.

\section{Future Horizons: Towards Algorithmic Invulnerability}
To expand upon the architecture established in this thesis, two rigorous operational vectors are planned for subsequent development, targeting absolute algorithmic validation and extreme infrastructure scalability.

\subsection{In Silico Validation Pipeline}
Establishing the absolute invulnerability of NECTA requires a strict computational proof of concept using synthetic benchmarks. Future developments will focus on generating highly controlled multivariate synthetic time series (Ground Truth), embedding predefined theoretical frequencies (30--90 Hz) and strictly orchestrated directional causal paths. By systematically introducing noise degradation (incorporating varying distributions of Pink and White noise), we will evaluate NECTA's resilience to adverse Signal-to-Noise Ratios (SNR). Spatial leakage and physical volume conduction artifacts will be modeled explicitly, allowing the extraction of exact precision metrics, such as Receiver Operating Characteristic (ROC) curves and Area Under the Curve (AUC) statistics, to mathematically benchmark true causal detection rates.

\subsection{HPC Architectural Stress Testing}
To decisively tackle the Big Data processing bottleneck inherent to multi-channel electrophysiology, the pipeline must undergo infrastructure-level architectural stress tests. Future work will benchmark standard sequential pandas-based data execution against parallel distributed orchestration managed via Dask. By tracking scalability across a multi-node high-performance computing environment using 1, 2, 4, 8, and 16 concurrent workers, we will map the acceleration profiles, memory footprint limits, and absolute latency boundaries of the framework. This benchmarking will ensure NECTA's scalability when executing time-resolved causal algorithms over massive, ultra-high-density recording sessions seamlessly.